\documentclass[12pt]{article}
\usepackage[utf8]{inputenc}
\usepackage[a4paper, margin=0.7in]{geometry} 
\usepackage{graphicx}
\usepackage{caption}
\usepackage{subcaption}
\captionsetup{skip=4pt}
\usepackage{booktabs}
\usepackage[table]{xcolor}
\usepackage{colortbl}
\usepackage[utf8]{inputenc}
\usepackage{longtable}
\usepackage{array}
\usepackage{geometry}
\usepackage{ragged2e}
\geometry{a4paper, margin=0.75in}
\usepackage{booktabs}
\usepackage{float}
\usepackage{fvextra}
\usepackage{listings}
\lstset{breaklines=true}

\title{Performing TLS Fingerprinting on
Network Namespaces}
\author{Samyak Gedam (252IS032)}
\date{March 2026}

\begin{document}

\maketitle


\noindent
\textbf{Aim:} \\

To analyze and compare TLS ClientHello fingerprints generated by different TLS clients(TLS v1.2 and v1.3) using packet captures and the Mercury fingerprinting framework. \\

\noindent
\textbf{Objective:} \\
\begin{itemize}
\item Capture TLS handshake traffic in an isolated Linux network namespace environment.
\item Generate TLS connections using different clients.
\item Extract TLS fingerprints using Mercury.
\item Compare fingerprints to observe differences in cipher suites and extensions. \\
\end{itemize}

\noindent
\textbf{Tools Used:} \\

Linux Network Namespaces, tcpdump, OpenSSL, curl, Cisco Mercury. \\

\section*{Network Namespaces Topology:}

The experiment was conducted in an isolated virtual network created using Linux network namespaces.

\begin{itemize}
\item red: TLS client namespace
\item router: packet capture namespace
\item blue: TLS server namespace \\
\end{itemize}


\begin{figure}
    \centering
    \includegraphics[width=\linewidth]{topology.png}
    \caption{Network Topology}
\end{figure}

\newpage
\noindent
\textbf{Packet Captures: } \\

The following PCAP files were generated.

\begin{itemize}
\item curl\_tls12.pcap
\item curl\_tls13.pcap
\item python\_ssl.pcap
\item openssl.pcap
\end{itemize}

\noindent
\textbf{Capturing TLS Fingerprint:} \\

Mercury was used to extract fingerprints.

\begin{lstlisting}[language=bash]
./src/mercury -r ../pcaps/<name>.pcap
\end{lstlisting}

\vspace{0.5cm}

\noindent
\textbf{Example output:}
\begin{verbatim}
tls/(0303)(130213031301...)
\end{verbatim}

\noindent
\textbf{Fingerprint format: }
\begin{verbatim}
tls/(TLS_VERSION)(CIPHER_SUITES)(EXTENSIONS)
\end{verbatim}


\noindent
\textbf{Screenshots: }\\
\begin{figure}
    \centering
    \includegraphics[width=\linewidth]{1.curl_tls12.png}
    \caption{Client - curl TLS1.2}
\end{figure}

\begin{figure}
    \centering
    \includegraphics[width=\linewidth]{2.curl_tls13.png}
    \caption{Client - curl TLS1.3}
\end{figure}

\begin{figure}
    \centering
    \includegraphics[width=\linewidth]{3.openssl.png}
    \caption{Client - OpenSSL}
\end{figure}

\begin{figure}
    \centering
    \includegraphics[width=\linewidth]{4.python_ssl.png}
    \caption{Client - Python SSL}
\end{figure}

\begin{figure}
    \centering
    \includegraphics[width=\linewidth]{F1.png}
\end{figure}
\begin{figure}
    \centering
    \includegraphics[width=\linewidth]{F2.png}
\end{figure}
\begin{figure}
    \centering
    \includegraphics[width=\linewidth]{F3.png}
\end{figure}

\clearpage
\section*{Results}

The extracted fingerprints from different TLS clients were compared.

\begin{table}[H]
\centering
\begin{tabular}{lll}
\toprule
Client & TLS Version & Observation \\
\midrule
curl TLS1.2 & TLS 1.2 & Large cipher suite list \\
curl TLS1.3 & TLS 1.3 & Modern AEAD cipher suites \\
Python SSL & TLS 1.3 & Different extension ordering \\
OpenSSL & TLS 1.2 & Minimal handshake structure \\
\bottomrule
\end{tabular}
\caption{Comparison of TLS fingerprints}
\end{table}

\section*{Discussion}

The results show that different TLS implementations produce distinct fingerprints even when connecting to the same server. The differences arise from:

\begin{itemize}
\item Cipher suite ordering
\item Supported TLS versions
\item TLS extensions
\item Client library implementation
\end{itemize}


\begin{table}[H]
\centering
\begin{tabular}{lll}
\toprule
Feature & TLS 1.2 & TLS 1.3 \\
\midrule
Cipher Suites & Large and diverse list & Short modern list \\
Typical Ciphers & c02c, c030, 009f, cca9 & 1301, 1302, 1303 \\
New Extensions & Not present & supported\_versions, key\_share \\
Handshake Structure & Multiple negotiation steps & Simplified and faster \\
Fingerprint Pattern & Longer and more variable & Shorter but distinct \\
\bottomrule
\end{tabular}
\caption{Comparison of TLS 1.2 and TLS 1.3 fingerprints}
\end{table}

\subsection*{Other TLS Fingerprinting Tools}

I identified several other tools besides Mercury. 
\textbf{JA3} is one of the most widely used methods and generates fingerprints by 
hashing fields from the TLS ClientHello message. \\
Some other Tools such as \textbf{Wireshark, Zeek, and Suricata} also support JA3 fingerprint extraction from packet captures. \\
But in comparison, \textbf{Mercury} provides more detailed fingerprints without hashing the handshake structure, which helped me for better analysis of TLS client behavior.

\section*{Conclusion}

Demonstrated how TLS fingerprinting can identify different TLS clients using handshake metadata. By capturing traffic and analyzing fingerprints with the Mercury framework, it was observed that each client produces a unique fingerprint due to variations in cipher suites and extensions.

\end{document}